% ============================================================
% 3.3 Định nghĩa các tình huống sử dụng (tổng quan)
% ============================================================

\section{Định nghĩa các tình huống sử dụng}
\label{sec:usecases-overview}

\subsection{Tác nhân hệ thống}

\begin{table}[htbp]
\centering
\caption{Các tác nhân trong module quản lý đơn hàng}
\label{tab:actors}
\small
\begin{tabular}{@{} c l p{8.5cm} @{}}
\toprule
\rowcolor{headerblue}
\thead{STT} & \thead{Tác nhân} & \thead{Vai trò} \\
\midrule
1 & Khách hàng (Buyer) & Mua sắm, giỏ hàng, đặt hàng, thanh toán, hoàn trả, tranh chấp \\
2 & Người bán (Seller) & Xác nhận/từ chối đơn, duyệt hoàn trả, dashboard phân tích \\
3 & Quản trị viên (Admin) & Giải quyết tranh chấp hoàn trả \\
\bottomrule
\end{tabular}
\end{table}

\subsection{Lược đồ use case tổng quan}

\begin{center}
\hspace*{-1cm} % <-- chỉnh giá trị âm tại đây
\begin{tikzpicture}[
  x=1.05cm, y=1.15cm, % scale nhẹ, tăng Y cho thoáng hàng
  uc/.style={
    ellipse, draw=headerblue, fill=rowgray, 
    text width=2.2cm, % giảm width cho gọn
    align=center, font=\footnotesize, 
    minimum height=0.8cm, inner sep=2.5pt
  },
  actor/.style={font=\footnotesize\bfseries},
  arr/.style={-, headerblue!60},
  inc/.style={-{Stealth[length=4pt]}, dashed, headerblue!60},
  ext/.style={-{Stealth[length=4pt]}, dotted, headerblue!60},
  lbl/.style={font=\scriptsize\itshape, headerblue!70!black},
  cnt/.style={font=\tiny, gray, anchor=north},
]

% System boundary (nới nhẹ)
\draw[headerblue, thick, rounded corners=8pt] (-0.5,0.8) rectangle (10.8,-11);

% Actors — Buyer
\node[actor] at (-2.3,-3) {Khách hàng};
\draw (-2.3,-3.5) circle(0.25);
\draw (-2.3,-3.75) -- (-2.3,-4.4);
\draw (-2.7,-4) -- (-1.9,-4);
\draw (-2.3,-4.4) -- (-2.7,-4.9);
\draw (-2.3,-4.4) -- (-1.9,-4.9);

% Seller
\node[actor] at (12.6,-3) {Người bán};
\draw (12.6,-3.5) circle(0.25);
\draw (12.6,-3.75) -- (12.6,-4.4);
\draw (12.2,-4) -- (13,-4);
\draw (12.6,-4.4) -- (12.2,-4.9);
\draw (12.6,-4.4) -- (13,-4.9);

% Admin
\node[actor] at (12.6,-9.3) {Quản trị viên};
\draw (12.6,-9.8) circle(0.25);
\draw (12.6,-10.05) -- (12.6,-10.7);
\draw (12.2,-10.3) -- (13,-10.3);
\draw (12.6,-10.7) -- (12.2,-11.2);
\draw (12.6,-10.7) -- (13,-11.2);

% ===== GRID GỌN & ĐỀU =====

% Row 1
\node[uc] (g1) at (1.8,0) {Giỏ hàng};
\node[cnt] at (g1.south) {UC01--03};

\node[uc] (g2) at (5.4,0) {Đặt hàng};
\node[cnt] at (g2.south) {UC04--06};

\node[uc] (g3) at (9,0) {Xử lý đơn};
\node[cnt] at (g3.south) {UC07--11};

% Row 2
\node[uc] (g4) at (1.8,-2.8) {Ví \& TT};
\node[cnt] at (g4.south) {UC12--13};

\node[uc] (g5) at (5.4,-2.8) {Theo dõi};
\node[cnt] at (g5.south) {UC14--16};

% Row 3
% Row 3 (fix spacing)
\node[uc] (g6) at (1.8,-5.8) {Hoàn trả};
\node[cnt] at (g6.south) {UC17--21};

\node[uc] (g7) at (6.6,-5.8) {Tranh chấp}; % <-- đẩy sang phải
\node[cnt] at (g7.south) {UC22--23};

% Row 4
\node[uc] (g9) at (1.8,-8.8) {Đánh giá};
\node[cnt] at (g9.south) {UC28--30};

\node[uc] (g8) at (5.4,-8.8) {Dashboard};
\node[cnt] at (g8.south) {UC24--27};

% ===== CONNECTIONS =====

% Buyer
\foreach \n in {g1,g2,g4,g5,g6,g9}
  \draw[arr] (-2.3,-4.2) -- (\n);
\draw[arr] (-2.3,-4.2) -- (g7);

% Seller
\foreach \n in {g3,g5,g6,g8}
  \draw[arr] (12.6,-4.2) -- (\n);
\draw[arr] (12.6,-4.2) -- (g7);

% Admin
\draw[arr] (12.6,-10) -- (g7);

% Relations
\draw[inc] (g2) -- (g4) node[lbl, midway, right] {$\ll$include$\gg$};
\draw[ext] (g7) -- (g6) node[lbl, midway, above] {$\ll$extend$\gg$};

\end{tikzpicture}
\captionof{figure}{Lược đồ use case tổng quan --- Module quản lý đơn hàng}
\label{fig:overall-usecase}
\end{center}

\subsection{Danh sách tình huống sử dụng}

\rowcolors{2}{white}{rowgray}
\begin{longtable}{@{} c p{7.5cm} p{4cm} @{}}
\caption{Danh sách tình huống sử dụng}\label{tab:usecase-list} \\
\toprule
\rowcolor{headerblue}
\thead{Mã} & \thead{Tên} & \thead{Actor} \\
\midrule
\endfirsthead
\toprule
\rowcolor{headerblue}
\thead{Mã} & \thead{Tên} & \thead{Actor} \\
\midrule
\endhead
\midrule
\multicolumn{3}{r}{\small\textit{Tiếp theo trang sau}} \\
\endfoot
\bottomrule
\endlastfoot
\rowcolor{white}\multicolumn{3}{@{}l}{\textit{\textbf{Giỏ hàng}}} \\
UC01 & Xem giỏ hàng & Buyer \\
UC02 & Cập nhật giỏ hàng & Buyer \\
UC03 & Xóa toàn bộ giỏ hàng & Buyer \\
\rowcolor{white}\multicolumn{3}{@{}l}{\textit{\textbf{Đặt hàng}}} \\
UC04 & Đặt hàng từ giỏ hàng & Buyer \\
UC05 & Mua ngay (Buy Now) & Buyer \\
UC06 & Hủy mục chờ xử lý & Buyer \\
\rowcolor{white}\multicolumn{3}{@{}l}{\textit{\textbf{Xử lý đơn (Seller)}}} \\
UC07 & Xem danh sách mục chờ & Seller \\
UC08 & Xem chi tiết mục chờ & Seller \\
UC09 & Xác nhận đơn hàng & Seller \\
UC10 & Từ chối đơn hàng & Seller \\
UC11 & Xem đơn đã xác nhận & Seller \\
\rowcolor{white}\multicolumn{3}{@{}l}{\textit{\textbf{Ví \& Thanh toán}}} \\
UC12 & Quản lý ví (xem số dư, lịch sử) & Buyer \\
UC13 & Xác nhận TT (webhook) & Hệ thống \\
\rowcolor{white}\multicolumn{3}{@{}l}{\textit{\textbf{Theo dõi}}} \\
UC14 & Xem đơn hàng (Buyer) & Buyer \\
UC15 & Xem đơn hàng (Seller) & Seller \\
UC16 & Xem DS đơn hàng & Buyer, Seller \\
\rowcolor{white}\multicolumn{3}{@{}l}{\textit{\textbf{Hoàn trả}}} \\
UC17 & Xem DS hoàn trả & Buyer \\
UC18 & Tạo yêu cầu hoàn trả & Buyer \\
UC19 & Cập nhật hoàn trả & Buyer \\
UC20 & Hủy hoàn trả & Buyer \\
UC21 & Duyệt hoàn trả & Seller \\
\rowcolor{white}\multicolumn{3}{@{}l}{\textit{\textbf{Tranh chấp}}} \\
UC22 & Mở tranh chấp & Buyer, Seller \\
UC23 & Giải quyết tranh chấp & Admin\textsuperscript{*} \\
\rowcolor{white}\multicolumn{3}{@{}l}{\textit{\textbf{Dashboard}}} \\
UC24 & Xem thống kê & Seller \\
UC25 & Xem xu hướng & Seller \\
UC26 & Xem hành động chờ & Seller \\
UC27 & Xem SP bán chạy & Seller \\
\rowcolor{white}\multicolumn{3}{@{}l}{\textit{\textbf{Đánh giá}}} \\
UC28 & Kiểm tra quyền đánh giá & Buyer \\
UC29 & Xem đơn có thể đánh giá & Buyer \\
UC30 & Xác thực đơn cho đánh giá & Buyer \\
\end{longtable}
\rowcolors{1}{}{}

\noindent\textsuperscript{*}UC23 (Giải quyết tranh chấp) yêu cầu vai trò Admin. Trong phạm vi đồ án, Admin giải quyết thủ công ngoài hệ thống; cơ sở dữ liệu đã dự phòng các trường \texttt{resolved\_by\_id}, \texttt{resolution\_note}, \texttt{date\_resolved} để hỗ trợ triển khai đầy đủ trong tương lai.

\subsection{Quan hệ và tương tác liên module}

\textbf{Generalization}: UC05 đặc biệt hóa UC04; UC14/UC15 tổng quát hóa ``Xem đơn hàng''. \textbf{$\ll$include$\gg$}: UC04 $\rightarrow$ đặt chỗ tồn kho + báo giá vận chuyển + thanh toán (bao gồm trừ ví); UC06/UC10 $\rightarrow$ giải phóng tồn kho + hoàn ví; UC09 $\rightarrow$ tạo vận đơn; UC21 $\rightarrow$ auto-refund. \textbf{$\ll$extend$\gg$}: UC05 mở rộng UC04; UC22 mở rộng UC21; UC13 mở rộng UC04 (webhook xác nhận thanh toán checkout).

Tương tác liên module qua Restate proxy: Catalog (tra cứu SKU), Inventory (tồn kho), Account (thông báo, địa chỉ, ví), Common (tài nguyên), Analytic (hành vi). Catalog gọi ngược Order để xác thực quyền đánh giá.
